\section*{Bab \ref{ch:g12:limcont} --- Limit dan Kekontinuan}

\begin{solution}{exo:g12:limcont:1}
Setelah dibagi $x^3$:
$\dfrac{2x^3 - x + 1}{5x^3 + x^2} = \dfrac{2 - 1/x^2 + 1/x^3}{5 + 1/x}
\to \dfrac{2}{5}$.

$x^2 + 3x - 1 = x^2(1 + 3/x - 1/x^2) \to +\infty$ ketika $x \to -\infty$, sebab
$x^2 \to +\infty$ dan isi kurungnya menuju $1$.

Ketika $x \to 2^+$, pembilangnya menuju $3 > 0$ dan penyebutnya menuju $0^+$,
sehingga $\dfrac{x+1}{x-2} \to +\infty$.

Untuk $x > 0$, $-\dfrac1x \leq \dfrac{\cos x}{x} \leq \dfrac1x$, sehingga
limitnya $0$ menurut teorema apit.
\end{solution}

\begin{solution}{exo:g12:limcont:2}
\emph{1.} $2x^2 - 3x + 1 = (x-1)(2x-1)$ (\omterm{def:g11:quad:discriminant}{akar} $x=1$ terlihat dari
$2 - 3 + 1 = 0$). Untuk $x \neq 1$, hapuslah $x - 1$: $f(x) = 2x - 1$.

\emph{2.} Jadi $\lim_{x\to1} f(x) = 1$. Menetapkan $f(1) = 1$ memperluas $f$
menjadi \omterm{def:g11:func:function}{fungsi} polinomial $x \mapsto 2x-1$ pada $\R$, yang \omterm{def:g12:limcont:continuity}{kontinu} di $1$.
\end{solution}

\begin{solution}{exo:g12:limcont:3}
Tulislah $f(x) = \dfrac{3x - 1}{x + 2} = \dfrac{3(x+2) - 7}{x+2}
= 3 - \dfrac{7}{x+2}$. Ketika $x \to \pm\infty$, $f(x) \to 3$: jadi garis
$y = 3$ adalah \omterm{def:g12:limcont:asymptote}{asimtot mendatar} di kedua takhingga. Ketika $x \to (-2)^\pm$,
$f(x) \to \mp\infty$: garis $x = -2$ adalah \omterm{def:g12:limcont:asymptote}{asimtot tegak}. Kurvanya berupa
\omterm{def:g11:func:reference}{hiperbola}, di bawah $y=3$ untuk $x > -2$ dan di atasnya untuk $x < -2$.
\end{solution}

\begin{solution}{exo:g12:limcont:4}
Kalikan dengan bentuk sekawannya:
\[
\sqrt{x^2 + x} - x = \frac{x^2 + x - x^2}{\sqrt{x^2+x} + x}
= \frac{x}{x\left(\sqrt{1 + 1/x} + 1\right)}
= \frac{1}{\sqrt{1 + 1/x} + 1}
\xrightarrow[x\to+\infty]{} \frac{1}{2}.
\]
Serupa dengan itu,
\[
\frac{\sqrt{1+x} - 1}{x} = \frac{(1+x) - 1}{x\left(\sqrt{1+x} + 1\right)}
= \frac{1}{\sqrt{1+x} + 1} \xrightarrow[x\to0]{} \frac{1}{2}.
\]
\end{solution}

\begin{solution}{exo:g12:limcont:5}
\emph{1.} Karena $\sin x \geq -1$, berlaku $f(x) \geq x - 1 \to +\infty$, dan
kita simpulkan lewat pembandingan. Namun $f'(x) = 1 + \cos x$ lenyap di setiap
kelipatan ganjil $\pi$, dan pada setiap \omterm{def:g10:numbers:interval}{interval} $\intco{A}{+\infty}$ \omterm{def:g11:func:function}{fungsi}
itu berselang-seling dengan \omterm{def:g10:numbers:interval}{interval} kenaikan yang tegas; \omterm{def:g11:func:function}{fungsi} itu naik tetapi
tidak tegas, dan yang lebih penting pertumbuhannya tidak diperlukan: pembandingannya sudah cukup.

\emph{2.} Untuk $x > 1$, bagilah dengan $x$:
\[
g(x) = \frac{1 + \frac{\sin x}{x}}{1 - \frac{\sin x}{x}},
\]
dan $\frac{\sin x}{x} \to 0$ menurut teorema apit
($\abs{\sin x} \leq 1$). Jadi $g(x) \to \frac{1+0}{1-0} = 1$.
\end{solution}

\begin{solution}{exo:g12:limcont:6}
$f(x) = x^3 - 3x + 1$ \omterm{def:g12:limcont:continuity}{kontinu}, dengan $f'(x) = 3x^2 - 3 = 3(x-1)(x+1)$:
jadi $f$ naik pada $\intoc{-\infty}{-1}$, turun pada $\intcc{-1}{1}$,
naik pada $\intco{1}{+\infty}$, dengan \omterm{def:g10:functions:extrema}{maksimum} lokal $f(-1) = 3 > 0$ dan
\omterm{def:g10:functions:extrema}{minimum} lokal $f(1) = -1 < 0$, $\lim_{-\infty} f = -\infty$,
$\lim_{+\infty} f = +\infty$.

Pada masing-masing ketiga \omterm{def:g10:numbers:interval}{interval} kemonotonannya, $f$ \omterm{def:g12:limcont:continuity}{kontinu}, tegas
\omterm{def:g12:seq:monotonic}{monoton}, dan $0$ terletak tegas di antara nilai atau limit di titik ujungnya,
sehingga teorema bijeksi memberi tepat satu nol per \omterm{def:g10:numbers:interval}{interval}: seluruhnya tiga
nol. Pemeriksaan tandanya: $f(-2) = -1 < 0 < f(-1) = 3$, jadi satu nol di
$\intoo{-2}{-1}$; $f(0) = 1 > 0 > f(1) = -1$, satu nol di $\intoo{0}{1}$;
$f(1) = -1 < 0 < f(2) = 3$, satu nol di $\intoo{1}{2}$.
\end{solution}

\begin{solution}{exo:g12:limcont:7}
Misalkan $g(x) = f(x) - x$, yang \omterm{def:g12:limcont:continuity}{kontinu} pada $\intcc{0}{1}$. Karena
$f(0) \in \intcc{0}{1}$, berlaku $g(0) = f(0) \geq 0$; karena
$f(1) \in \intcc{0}{1}$, berlaku $g(1) = f(1) - 1 \leq 0$. Menurut teorema
nilai antara, $g$ lenyap di suatu $c \in \intcc{0}{1}$, yaitu $f(c) = c$.
\end{solution}

\begin{solution}{exo:g12:limcont:8}
Misalkan $u(t)$ dan $d(t)$ kedudukannya (diukur sebagai jarak dari kaki
jalurnya) pada waktu $t \in \intcc{8}{20}$ ketika mendaki dan ketika menuruni
jalur. Keduanya \omterm{def:g12:limcont:continuity}{kontinu}, dengan $u(8) = 0$, $u(20) = L$ (panjang jalurnya),
$d(8) = L$, $d(20) = 0$. \omterm{def:g11:func:function}{Fungsi} $g = u - d$ \omterm{def:g12:limcont:continuity}{kontinu} dengan
$g(8) = -L \leq 0$ dan $g(20) = L \geq 0$, sehingga $g(c) = 0$ untuk suatu
waktu $c$ menurut teorema nilai antara: pada waktu $c$, kedudukannya pada
kedua hari itu berimpit.
\end{solution}

\begin{solution}{exo:g12:limcont:9}
\emph{1.} $f'(x) = 5x^4 + 3x^2 + 1 > 0$, sehingga $f$ tegas naik pada
$\R$; dengan $\lim_{\pm\infty} f = \pm\infty$ dan \omterm{def:g12:limcont:continuity}{kekontinuannya}, teorema
bijeksi memberi satu nol real tunggal $\alpha$. Karena $f(0) = -1 < 0$ dan
$f(1) = 2 > 0$, diperoleh $\alpha \in \intoo{0}{1}$.

\emph{2.} Pertahankan sebuah \omterm{def:g10:numbers:interval}{interval} $\intcc{a_n}{b_n}$ dengan
$f(a_n) < 0 < f(b_n)$, dimulai dari $\intcc{0}{1}$; pada setiap langkahnya
hitunglah $f$ di \omterm{prop:g10:coordgeom:midpoint}{titik tengah} $m$ lalu simpan separuh \omterm{def:g10:numbers:interval}{interval} tempat
pergantian tandanya bertahan. Setelah $n$ langkah, \omterm{def:g10:numbers:interval}{intervalnya} berpanjang
$2^{-n}$, sehingga \omterm{prop:g10:coordgeom:midpoint}{titik tengahnya} menghampiri $\alpha$ dengan ketelitian
$2^{-(n+1)}$. Kita memerlukan $2^{-(n+1)} \leq 10^{-6}$, yaitu
$n + 1 \geq \frac{6}{\log_{10} 2} \approx 19.93$: jadi $n = 20$ iterasi sudah cukup.
\end{solution}

\begin{solution}{pb:g12:limcont:1}
\textbf{1.} Suku yang menguasai: $\frac{3x^2}{x^2} \to 3$;
$\frac{2x}{x^2} = \frac2x \to 0$;
$x^3 - x^2 = x^3\left(1 - \frac1x\right) \to +\infty$.

\textbf{2.} $\lim_{x \to \pm\infty} f = 3$: jadi \omterm{def:g12:limcont:asymptote}{asimtot
mendatarnya} $y = 3$; $\lim_{x \to (-2)^\pm} f = \mp\infty$:
\omterm{def:g12:limcont:asymptote}{asimtot tegaknya} $x = -2$.

\textbf{3.} $\frac{(x-3)(x+3)}{x - 3} = x + 3 \to 6$. Lalu
$\frac{\sqrt{x+1} - 1}{x} = \frac{x}{x(\sqrt{x+1} + 1)} =
\frac{1}{\sqrt{x+1} + 1} \to \frac12$.

\textbf{4.} $\abs{x \sin\frac1x} \leq \abs x \to 0$, sehingga
limitnya $0$. Lalu
$\abs{\frac{\sin x}{x}} \leq \frac1x \to 0$: limitnya $0$.

\textbf{5.} $\frac{\sqrt{4x^2 + x}}{x} =
\sqrt{4 + \frac1x} \to \sqrt4 = 2$ (lewat \omterm{def:g12:limcont:continuity}{kekontinuan} akar
kuadrat dan komposisi).

\textbf{6.} $f(x) = x^3 + x$ \omterm{def:g12:limcont:continuity}{kontinu} dengan
$f(1) = 2 < 5 < 10 = f(2)$: menurut teorema nilai antara, $f(c) = 5$ untuk
suatu $c \in \intoo{1}{2}$.

\textbf{7.} $f'(x) = 3x^2 + 1 > 0$: jadi tegas naik pada
$\R$, \omterm{def:g12:limcont:continuity}{kontinu}, dengan limit $\mp\infty$: menurut teorema bijeksi,
setiap nilai --- khususnya $5$ --- dicapai
tepat satu kali.

\textbf{8.} $p(x) = x^3\left(1 + \frac ax + \frac{b}{x^2} +
\frac{c}{x^3}\right)$: ketika $x \to +\infty$, $p \to +\infty$; ketika
$x \to -\infty$, $p \to -\infty$. Jadi pada suatu \omterm{def:g10:numbers:interval}{interval} yang sangat besar
$\intcc{-M}{M}$ berlaku $p(-M) < 0 < p(M)$: teorema nilai antara menyerahkan
sebuah \omterm{def:g11:quad:discriminant}{akar}. Alasan suku yang menguasai itu berlaku kata demi kata untuk
sebarang derajat ganjil --- pangkat ganjil mempertahankan tanda yang
berlawanan di kedua takhingga.

\textbf{9.} Membagi dua sebanyak $n$ kali menyisakan \omterm{def:g10:numbers:interval}{interval} sepanjang
$2^{-n}$; $2^{-20} \approx 9.5 \times 10^{-7} < 10^{-6}$:
jadi dua puluh langkah. Newton melipatduakan angka yang benar setiap langkah:
dari satu angka yang benar, sekitar $3$--$4$ langkah sudah mencapai enam ---
si kura-kura dan si anjing pemburu dalam pencarian \omterm{def:g11:quad:discriminant}{akar}.

\textbf{10.} $g$ tidak \omterm{def:g12:limcont:continuity}{kontinu} pada $\intcc{-1}{1}$: \omterm{def:g11:func:function}{fungsi} itu
meledak di $0$, yang termuat dalam \omterm{def:g10:numbers:interval}{intervalnya}. Kesimpulan teorema nilai
antara menguap bersama hipotesisnya --- periksalah selalu
\omterm{def:g12:limcont:continuity}{kekontinuannya} \emph{pada seluruh \omterm{def:g10:numbers:interval}{intervalnya}}, bukan hanya di
titik ujungnya.

\textbf{11.} $g(x) = f(x) - x$ \omterm{def:g12:limcont:continuity}{kontinu};
$g(0) = f(0) \geq 0$ (nilainya terletak di $\intcc{0}{1}$) dan
$g(1) = f(1) - 1 \leq 0$. Menurut teorema nilai antara, $g(c) = 0$ untuk suatu
$c$: yaitu $f(c) = c$.

\textbf{12.} Dengan mengidealkan wilayahnya sebagai ruas
$\intcc{0}{1}$: $f$ mengirim setiap tempat ke tempat yang tergambar
tepat di bawahnya pada \omterm{def:g10:functions:function}{peta} yang teremas; meremas \omterm{def:g10:functions:function}{peta} itu tetap membuatnya
berada di atas wilayahnya, sehingga $f$ mendarat di $\intcc{0}{1}$, dan
\omterm{def:g11:func:function}{fungsi} itu \omterm{def:g12:limcont:continuity}{kontinu} (kertasnya tidak robek). \omterm{pb:g10:functions:1}{Titik tetap} $c$ adalah
tempat yang terletak persis di atas \omterm{def:g10:functions:function}{petanya} sendiri --- teorema Brouwer
menangani \omterm{def:g10:functions:function}{peta} dua dimensi yang jujur.

\textbf{13.} $g(0) + g(\pi) = \left(T(0) - T(\pi)\right) +
\left(T(\pi) - T(2\pi)\right) = T(0) - T(2\pi) = 0$. Jadi
$g(0)$ dan $g(\pi)$ saling berlawanan (atau keduanya nol): bagaimanapun juga
$g$, yang \omterm{def:g12:limcont:continuity}{kontinu}, lenyap di suatu $t^* \in \intcc{0}{\pi}$:
$T(t^*) = T(t^* + \pi)$ --- jadi di suatu tempat pada khatulistiwa, dua
titik antipoda bersuhu sama, pada setiap saat.

\textbf{14.} Memutar mejanya seperempat putaran mempertukarkan
peran kedua diagonalnya: sepasang kaki yang menekan tanah dan sepasang kaki
yang (salah satunya) melayang bertukar keadaan, sehingga tinggi
layangnya --- yang dihitung bertanda sebagai ``celah diagonal yang
diistimewakan'' --- bertanda berlawanan pada $\theta$ dan
$\theta + \frac\pi2$. Tanah yang \omterm{def:g12:limcont:continuity}{kontinu} membuat $h$ \omterm{def:g12:limcont:continuity}{kontinu}, dan teorema
nilai antara menyerahkan sebuah sudut $\theta^*$ dengan $h(\theta^*) = 0$:
keempat kakinya menapak. (Meja persegi panjang atau tanah berundak tidak
seberuntung itu.)

\textbf{15.} Polanya: bangunlah \emph{selisih} kedua besaran
yang ingin kamu buat berimpit ($f(x) - x$;
$T(t) - T(t + \pi)$; kedudukan mendaki dikurangi kedudukan menuruni bagi
sang pendaki; celah bertanda bagi mejanya); periksalah bahwa selisih itu
mengambil nilai dari kedua tanda (sering lewat kesimetrian ujungnya); lalu
panggillah teorema nilai antara pada \omterm{def:g11:func:function}{fungsi} yang \omterm{def:g12:limcont:continuity}{kontinu}. Keberimpitannya terpaksa.

\textbf{16.} Pertanyaan yang terbuka: \emph{berapa banyak} penyelesaiannya
(dijawab oleh kemonotonan yang tegas dan teorema bijeksi,
pertanyaan 7) dan \emph{di mana} letaknya (dijawab oleh \omterm{thm:g12:limcont:ivt}{dikotomi}
atau Newton, pertanyaan 9). Teorema nilai antara adalah peramal keberadaan yang murni.

\textbf{17.} $3 - \frac{7}{x + 2}$: \omterm{def:g12:limcont:asymptote}{asimtot tegaknya}
$x = -2$, \omterm{def:g12:limcont:asymptote}{asimtot mendatarnya} $y = 3$; pada masing-masing sisi $-2$ \omterm{def:g11:func:function}{fungsi}
itu naik ($-\frac{7}{x+2}$ naik di tempat yang terdefinisi); kurvanya adalah
\omterm{def:g11:func:reference}{hiperbola} acuan $-\frac7x$ yang tergeser $2$ ke kiri dan $3$ ke atas ---
dua cabang yang memeluk perpotongan \omterm{def:g12:limcont:asymptote}{asimtotnya}.

\textbf{18.} Untuk $x \neq 1$: $f(x) = x + 1$, sehingga
$\lim_{x \to 1} f = 2$: menetapkan $f(1) = 2$ memulihkan
\omterm{def:g12:limcont:continuity}{kekontinuannya}. Ketakkontinuannya berupa tusukan, bukan lompatan ---
isilah satu titik yang hilang itu dan kurvanya sembuh:
\emph{terhapuskan}.

\textbf{19.} \omterm{def:g12:limcont:continuity}{Kontinu} pada setiap \omterm{def:g10:numbers:interval}{interval} antara dua \omterm{def:g10:numbers:sets}{bilangan bulat}
berurutan; di setiap \omterm{def:g10:numbers:sets}{bilangan bulat} $n$ \omterm{def:g11:func:function}{fungsi} itu melompat dari $n - 1$ ke $n$.
Dengan $a = 0.9$, $b = 1.1$: $\lfloor a \rfloor = 0 < \frac12 <
1 = \lfloor b \rfloor$, namun $\lfloor x \rfloor = \frac12$
tidak berpenyelesaian --- lompatannya berteleportasi melewati $\frac12$,
dan justru itulah yang dilarang teorema nilai antara bagi \omterm{def:g11:func:function}{fungsi} \omterm{def:g12:limcont:continuity}{kontinu}.

\textbf{20.} Wajah lokalnya: nilai $=$ limit di setiap titik ---
tusukan yang terhapuskan (pertanyaan 18) melanggarnya tetapi masih dapat
diobati, sedangkan lompatan \omterm{def:g11:func:function}{fungsi} lantai (pertanyaan 19) tidak. Wajah
globalnya: pada sebuah \omterm{def:g10:numbers:interval}{interval}, semua nilai antaranya disinggahi --- surat
keterangan \omterm{def:g11:quad:discriminant}{akar} (6--8), \omterm{pb:g10:functions:1}{titik tetapnya}, khatulistiwanya, mejanya
dan pendakinya semuanya tinggal di sini. Nama resmi semboyannya:
\emph{teorema nilai antara}.

\end{solution}
